\documentclass[11pt,a4paper]{article}

\usepackage[utf8]{inputenc}
\usepackage[T1]{fontenc}
\usepackage[french]{babel}
\usepackage{lmodern}
\usepackage{geometry}
\usepackage{amsmath,amssymb,amsfonts}
\usepackage{bm}
\usepackage{graphicx}
\usepackage{booktabs}
\usepackage{hyperref}
\usepackage{xcolor}
\usepackage[most]{tcolorbox}

\geometry{margin=2.5cm}

\title{
\textbf{Source tensorielle émergente de la matière}\\[0.4em]
\large De l'excitation locale \(\delta W^{\rm loc}\) à une source matière BuP candidate
}

\author{
Farid Hamdad\\
\small Bottom-Up Quantum Gravity Program
}

\date{2026}

\begin{document}

\maketitle

\begin{abstract}

Dans le programme Bottom-Up Quantum Gravity, la matière ne doit pas être
introduite comme une source externe ajoutée à la géométrie. Elle doit émerger
du réseau d'intrication lui-même. Le présent papier construit et teste une
source tensorielle candidate à partir d'une perturbation locale du réseau
d'information mutuelle,
\[
W_{ij}=W_{ij}^{(0)}+\delta W_{ij}^{\rm loc}.
\]

Nous définissons d'abord une densité d'énergie informationnelle,
\[
T_{00}^{\rm matter,cand}(i)
=
\frac12
\sum_j(\delta W_{ij}^{\rm loc})^2.
\]
Cette quantité est positive, correctement normalisée et localisée. Cependant,
elle ne suffit pas à prédire la réponse géométrique locale. Nous introduisons
alors successivement une trace spatiale isotrope \(T_{aa}\), un terme de
contrainte interne \(T_{\rm grad}\), puis une composante de flux \(T_{0a}\).

La meilleure source candidate testée est
\[
S_{\rm flux}
=
T_{00}
-\frac12T_{aa}
+
\frac12T_{\rm grad}
+
\|T_{0a}\|.
\]
Elle atteint
\[
\rho_{\rm Spearman}(S_{\rm flux},|\delta R|)
=
0.741,
\qquad
p=1.84\times10^{-4}.
\]

Nous montrons également que la source matière candidate n'est pas conservée
isolément. Cette non-conservation n'est pas interprétée comme un échec, mais
comme une signature d'échange entre matière émergente et fond d'intrication.
Un flux d'intrication construit à partir de
\[
\phi_{\rm ent}
=
\delta d_s+\delta R+\frac12\Delta\delta R
\]
réduit le résidu de divergence d'un facteur \(2.63\). Ainsi, Paper~8 construit
une source tensorielle matière candidate fortement corrélée à la courbure et
partiellement compensée par le champ d'intrication.

\end{abstract}

\tableofcontents

\section{Introduction}

Dans la formulation effective de BuP, la géométrie émergente est reconstruite à
partir du réseau d'information mutuelle
\[
W_{ij}=I(i:j).
\]
Paper~5 a introduit une équation d'Einstein effective contenant un tenseur de
matière,
\[
G_{\mu\nu}[g^{\rm ent}]
=
8\pi G_{\rm eff}(d_s)T_{\mu\nu}^{\rm matter}
+
8\pi G T_{\mu\nu}^{\rm ent}[d_s].
\]

Paper~6 a proposé une équation pure BuP,
\[
\mathcal{F}_{ij}[W]=0,
\]
et Paper~7 a étudié la dynamique variationnelle des excitations localisées
\(\delta W^{\rm loc}\).

Le présent papier franchit l'étape suivante :
\[
\delta W^{\rm loc}
\longrightarrow
T_{\mu\nu}^{\rm matter,cand}.
\]

Nous insistons sur le terme \emph{candidate}. La conservation covariante
complète n'est pas démontrée d'emblée. Nous construisons une source tensorielle
effective, nous testons sa positivité, sa localisation, sa corrélation avec la
courbure et son échange avec le fond d'intrication.

\section{Excitation locale du réseau}

On considère un réseau de référence \(W^{(0)}\), puis une perturbation locale :
\[
W_{ij}
=
W_{ij}^{(0)}
+
\delta W_{ij}^{\rm loc}.
\]

Dans les simulations, l'excitation est injectée sous forme gaussienne en
distance d'intrication :
\[
\delta W_{ij}^{\rm loc}
=
A f_i f_j,
\]
avec
\[
f_i
=
\exp\left(
-\frac{d_{\rm ent}^2(i,i_0)}{2\sigma^2}
\right).
\]

\section{Densité d'énergie informationnelle}

La densité minimale est définie par
\[
T_{00}^{\rm matter,cand}(i)
=
\frac12
\sum_j
(\delta W_{ij}^{\rm loc})^2.
\]

Le facteur \(1/2\) évite le double comptage des arêtes :
\[
\sum_iT_{00}^{\rm matter,cand}(i)
=
\sum_{i<j}
(\delta W_{ij}^{\rm loc})^2.
\]

Cette densité est positive par construction :
\[
T_{00}^{\rm matter,cand}(i)\ge0.
\]

\section{Composantes tensorielles}

À partir d'un embedding de la géométrie émergente, on définit
\[
v_{ij}^a
=
\frac{x_j^a-x_i^a}{d_{ij}^{\rm ent}+\epsilon}.
\]

\subsection{Trace spatiale isotrope}

\[
T_{aa}(i)
=
\frac12
\sum_j
(\delta W_{ij}^{\rm loc})^2
\|v_{ij}\|^2.
\]

\subsection{Contrainte interne}

\[
T_{\rm grad}(i)
=
\sum_j
A_{ij}
\left(
T_{00}(i)-T_{00}(j)
\right)^2.
\]

\subsection{Flux}

Le flux informationnel est
\[
\vec f(i)
=
\sum_j
(\delta W_{ij}^{\rm loc})^2
\vec v_{ij},
\]
et sa norme est
\[
F(i)=\|\vec f(i)\|.
\]

Ce terme représente la composante \(T_{0a}\) de la source tensorielle candidate.

\section{Sources testées}

Nous testons une hiérarchie de sources.

\subsection{Densité seule}

\[
S_0(i)=T_{00}(i).
\]

\subsection{Source fluide}

\[
S_{\rm fluid}(i)
=
T_{00}(i)
+
\omega T_{aa}(i).
\]

Le meilleur poids effectif trouvé est :
\[
\omega=-0.5.
\]

\subsection{Source avec contrainte interne}

\[
S_{\rm full}(i)
=
T_{00}(i)
-\frac12T_{aa}(i)
+
\chi T_{\rm grad}(i),
\]
avec meilleur poids normalisé :
\[
\chi=0.5.
\]

\subsection{Source avec flux}

\[
S_{\rm flux}(i)
=
T_{00}(i)
-\frac12T_{aa}(i)
+
\frac12T_{\rm grad}(i)
+
\psi F(i).
\]

Le meilleur poids normalisé est :
\[
\psi=1.0.
\]

Ainsi, la meilleure source testée est :
\[
\boxed{
S_{\rm flux}
=
T_{00}
-\frac12T_{aa}
+
\frac12T_{\rm grad}
+
\|T_{0a}\|.
}
\]

\section{Protocole numérique}

Les tests principaux utilisent :
\[
N=20,
\qquad
\lambda=0.57,
\qquad
k=5,
\qquad
A=0.15.
\]

La réponse géométrique est :
\[
\delta R_i
=
R_i[W^{(0)}+\delta W]
-
R_i[W^{(0)}].
\]

Le critère principal est :
\[
\rho_{\rm Spearman}
\left(
S(i),
|\delta R_i|
\right).
\]

\section{Résultats : positivité et localisation}

Le test de positivité donne :
\[
T_{00}^{\rm matter,cand}(i)>0
\]
pour tous les nœuds.

La normalisation est vérifiée à précision machine :
\[
\sum_iT_{00}(i)
=
\sum_{i<j}(\delta W_{ij})^2.
\]

Le scan en largeur \(\sigma\) montre que \(T_{00}\) est radialement décroissant
avec la distance d'intrication au centre. Pour \(\sigma=0.01\), le participation
ratio est proche de \(2\), indiquant une excitation compacte. Pour des valeurs
plus grandes de \(\sigma\), l'excitation devient plus étendue tout en restant
localisée géométriquement.

\section{Résultats : corrélation avec la courbure}

La densité seule est insuffisante :
\[
T_{00}
\quad\Rightarrow\quad
\rho_{\rm Spearman}\simeq0.195,
\qquad
p\simeq0.409.
\]

L'ajout de la trace spatiale améliore le résultat :
\[
T_{00}-\frac12T_{aa}
\quad\Rightarrow\quad
\rho_{\rm Spearman}\simeq0.526,
\qquad
p\simeq0.017.
\]

L'ajout de la contrainte interne donne :
\[
T_{00}-\frac12T_{aa}+\frac12T_{\rm grad}
\quad\Rightarrow\quad
\rho_{\rm Spearman}\simeq0.612,
\qquad
p\simeq0.00413.
\]

Enfin, l'ajout du flux donne :
\[
T_{00}-\frac12T_{aa}+\frac12T_{\rm grad}+\|T_{0a}\|
\quad\Rightarrow\quad
\rho_{\rm Spearman}\simeq0.741,
\qquad
p\simeq1.84\times10^{-4}.
\]

\begin{table}[htbp]
\centering
\begin{tabular}{l c c c}
\toprule
Source testée & Paramètres & Spearman avec \(|\delta R|\) & \(p\)-value \\
\midrule
\(T_{00}\) & \(\sigma=0.05\) & \(0.195\) & \(0.409\) \\
\(T_{00}-\frac12T_{aa}\) & \(\sigma=0.15\) & \(0.526\) & \(0.017\) \\
\(T_{00}-\frac12T_{aa}+\frac12T_{\rm grad}\) & \(\sigma=0.15\) & \(0.612\) & \(0.00413\) \\
\(T_{00}-\frac12T_{aa}+\frac12T_{\rm grad}+\|T_{0a}\|\) & \(\sigma=0.15\) & \(0.741\) & \(1.84\times10^{-4}\) \\
\bottomrule
\end{tabular}
\caption{
Amélioration progressive de la source tensorielle candidate. La densité seule
est insuffisante ; les composantes spatiales, internes et de flux améliorent
fortement la corrélation avec la réponse de courbure.
}
\end{table}

\section{Conservation et échange matière--intrication}

Un tenseur énergie-impulsion physique conservé vérifie :
\[
\nabla_\mu T^{\mu\nu}=0.
\]

Cependant, dans BuP, la matière n'est pas un secteur fondamental isolé. Elle est
une excitation du réseau d'intrication. Il n'est donc pas nécessaire que
\[
\nabla_\mu T_{\rm matter}^{\mu\nu}=0
\]
soit vérifiée séparément.

L'équation effective complète impose plutôt :
\[
\nabla^\mu
\left[
G_{\rm eff}(d_s)T_{\mu\nu}^{\rm matter}
+
G T_{\mu\nu}^{\rm ent}
\right]
=
0.
\]

Nous interprétons donc la non-conservation isolée de la source matière candidate
comme un courant d'échange :
\[
\nabla_\mu T_{\rm matter}^{\mu\nu}
=
-J_{\rm exchange}^{\nu}.
\]

\section{Test de compensation totale}

Nous construisons un flux d'intrication à partir d'un champ
entropico-géométrique :
\[
\phi_{\rm ent}
=
\delta d_s
+
\beta\,\delta R
+
\gamma\,\Delta\delta R.
\]

Le flux associé est :
\[
\vec f_{\rm ent}(i)
=
-\sum_j A_{ij}
\left(
\phi_{\rm ent}(j)-\phi_{\rm ent}(i)
\right)
\vec v_{ij}.
\]

Le but est de réduire :
\[
\left\|
\mathrm{div}
\left(
\vec f_{\rm matter}
+
\alpha\vec f_{\rm ent}
\right)
\right\|_2.
\]

Le meilleur résultat est obtenu pour :
\[
\sigma=0.01,
\qquad
\beta=1.0,
\qquad
\gamma=0.5,
\]
soit :
\[
\phi_{\rm ent}
=
\delta d_s
+
\delta R
+
\frac12\Delta\delta R.
\]

On obtient :
\[
\|\mathrm{div}\,\vec f_{\rm matter}\|_2
=
1.62\times10^{-5},
\]
et :
\[
\|\mathrm{div}\,\vec f_{\rm total}\|_2
=
6.18\times10^{-6}.
\]

Le résidu relatif est :
\[
\frac{
\|\mathrm{div}\,\vec f_{\rm total}\|_2
}{
\|\mathrm{div}\,\vec f_{\rm matter}\|_2
}
=
0.380.
\]

La divergence est donc réduite d'un facteur :
\[
2.63.
\]

\begin{table}[htbp]
\centering
\begin{tabular}{l c c c}
\toprule
Test & Meilleur cas & Résidu relatif \(L^2\) & Interprétation \\
\midrule
Matière seule & -- & \(1.000\) & non conservée isolément \\
Compensation \(\delta d_s\) & \(\sigma=0.15\) & \(0.668\) & faible compensation \\
Compensation \(\delta d_s+\beta\delta R+\gamma\Delta\delta R\) & \(\sigma=0.01\) & \(0.380\) & compensation partielle sérieuse \\
\bottomrule
\end{tabular}
\caption{
Test de conservation totale matière--intrication. La source matière candidate
n'est pas conservée isolément, mais un flux d'intrication géométrique réduit la
divergence d'un facteur \(2.63\).
}
\end{table}

\section{Discussion}

Le résultat principal est double.

Premièrement, \(\delta W^{\rm loc}\) permet de construire une source tensorielle
candidate fortement corrélée à la réponse de courbure locale. La densité seule
ne suffit pas ; la pression isotrope, la contrainte interne et le flux sont
nécessaires.

Deuxièmement, la source matière candidate n'est pas conservée isolément. Cette
non-conservation devient une signature d'émergence : la matière échange avec le
fond d'intrication. La compensation partielle par un flux construit à partir de
\(\delta d_s\), \(\delta R\) et \(\Delta\delta R\) soutient cette interprétation.

\section{Limites}

La conservation totale n'est pas encore démontrée exactement. Le tenseur
d'intrication \(T_{\mu\nu}^{\rm ent}\) n'est représenté ici que par un proxy de
flux géométrique. Les coefficients numériques obtenus sont effectifs et
dépendent du protocole de normalisation.

Il faudra tester d'autres valeurs de \(N\), \(\lambda\), \(k\), amplitudes et
formes d'excitation.

\section{Conclusion}

Paper~8 construit une source tensorielle matière candidate à partir d'une
excitation locale du réseau d'intrication :
\[
\delta W^{\rm loc}
\longrightarrow
T_{\mu\nu}^{\rm matter,cand}.
\]

La meilleure source testée est :
\[
\boxed{
S_{\rm flux}
=
T_{00}
-\frac12T_{aa}
+
\frac12T_{\rm grad}
+
\|T_{0a}\|.
}
\]

Elle atteint :
\[
\rho_{\rm Spearman}(S_{\rm flux},|\delta R|)
=
0.741,
\qquad
p=1.84\times10^{-4}.
\]

La source n'est pas conservée isolément, mais sa divergence est partiellement
compensée par un flux d'intrication. Ce résultat transforme une faiblesse
apparente en signature d'émergence : la matière BuP n'est pas un secteur fermé,
mais une excitation échangeant avec le fond géométrique informationnel.

\[
\boxed{
\nabla_\mu
\left(
T_{\rm matter}^{\mu\nu}
+
T_{\rm ent}^{\mu\nu}
\right)
\simeq0
}
\]

reste l'objectif central de la suite du programme.

\begin{thebibliography}{9}

\bibitem{hamdad-paper5}
F.~Hamdad,
\textit{Équation d'Einstein effective émergente dans la Gravité Quantique Bottom-Up},
BuP Paper 5, 2026.

\bibitem{hamdad-paper6}
F.~Hamdad,
\textit{Équation pure BuP : de l'intrication à l'équation d'Einstein effective},
BuP Paper 6, 2026.

\bibitem{hamdad-paper7}
F.~Hamdad,
\textit{Dynamique variationnelle des excitations BuP},
BuP Paper 7, 2026.

\bibitem{ollivier}
Y.~Ollivier,
\textit{Ricci curvature of Markov chains on metric spaces},
Journal of Functional Analysis, 256, 810--864, 2009.

\end{thebibliography}

\end{document}
